\documentclass[11pt,a4paper]{article}
\usepackage[utf8]{inputenc}
\usepackage[T1]{fontenc}
\usepackage[brazilian]{babel}
\usepackage{geometry}
\usepackage{longtable}
\usepackage{booktabs}
\usepackage{array}
\usepackage{titling}

\geometry{margin=2.2cm}

\title{Verificação e Validação de Software \\ \large Trabalho 1: Casos de Teste da Tarifa de Estacionamento}
\author{}
\date{}

\begin{document}
\maketitle
\vspace{-2.5em}

\section*{Introdução}

Este documento apresenta o conjunto de casos de teste projetados para a classe
\texttt{CalculadoraTarifa}, responsável pelo cálculo do valor a ser pago pelo
ticket de estacionamento. Os casos foram derivados por particionamento em
classes de equivalência e por análise de valor limite, cobrindo as
variáveis de entrada: horário de entrada, horário de saída, duração de
permanência, condição de pernoite e status de cliente VIP.

\section*{Classes de equivalência identificadas}

\begin{itemize}
  \item \textbf{Horário de entrada}: válido no intervalo 08:00 a 23:59; inválido no
        intervalo 00:00 a 07:59.
  \item \textbf{Horário de saída}: válido no intervalo 08:00 a 01:59 (dia seguinte);
        inválido no intervalo 02:00 a 07:59.
  \item \textbf{Duração de permanência (sem pernoite)}: cortesia (até 20 minutos),
        tarifa fixa (21 a 60 minutos), tarifa fixa mais horas adicionais (acima de 60
        minutos).
  \item \textbf{Condição de pernoite}: saída antes das 08:00 do dia seguinte à
        entrada (sem pernoite) ou a partir das 08:00 do dia seguinte, inclusive
        (com pernoite, contando uma tarifa de pernoite por cada janela de 24 horas
        excedida).
  \item \textbf{Cliente}: VIP (desconto de 50\% sobre o valor final) ou não VIP.
  \item \textbf{Ordem cronológica}: saída posterior à entrada (válida) ou saída
        anterior ou igual à entrada (inválida).
\end{itemize}

\section*{Casos de teste: cálculo do valor da tarifa}

Todos os casos abaixo assumem entrada e saída válidas quanto ao horário. A coluna
\emph{Valor esperado} é o retorno do método \texttt{calcular}.

\begin{longtable}{@{}p{1.4cm}p{2.6cm}p{2.9cm}p{0.8cm}p{2.7cm}p{2.3cm}@{}}
\toprule
\textbf{ID} & \textbf{Entrada} & \textbf{Saída} & \textbf{VIP} & \textbf{Técnica / limite} & \textbf{Valor esperado} \\
\midrule
\endhead
CT01 & 10/03/2026 08:00 & 10/03/2026 08:19 & Não & Partição: cortesia (19 min) & R\$ 0,00 \\
CT02 & 10/03/2026 08:00 & 10/03/2026 08:20 & Não & Limite: fim da cortesia (20 min) & R\$ 0,00 \\
CT03 & 10/03/2026 08:00 & 10/03/2026 08:21 & Não & Limite: início da tarifa fixa (21 min) & R\$ 15,00 \\
CT04 & 10/03/2026 08:00 & 10/03/2026 08:59 & Não & Partição: tarifa fixa (59 min) & R\$ 15,00 \\
CT05 & 10/03/2026 08:00 & 10/03/2026 09:00 & Não & Limite: fim da tarifa fixa (60 min) & R\$ 15,00 \\
CT06 & 10/03/2026 08:00 & 10/03/2026 09:01 & Não & Limite: início da 1\textsuperscript{a} hora adicional (61 min) & R\$ 20,00 \\
CT07 & 10/03/2026 08:00 & 10/03/2026 10:00 & Não & Limite: fim da 1\textsuperscript{a} hora adicional (120 min) & R\$ 20,00 \\
CT08 & 10/03/2026 08:00 & 10/03/2026 10:01 & Não & Limite: início da 2\textsuperscript{a} hora adicional (121 min) & R\$ 25,00 \\
CT09 & 10/03/2026 08:00 & 10/03/2026 12:00 & Não & Partição: permanência de 4 horas & R\$ 30,00 \\
CT10 & 10/03/2026 08:00 & 10/03/2026 08:15 & Não & Limite: início do horário de entrada permitido (08:00) & R\$ 0,00 \\
CT11 & 10/03/2026 23:59 & 11/03/2026 00:30 & Não & Limite: fim do horário de entrada permitido (23:59) & R\$ 15,00 \\
CT12 & 10/03/2026 23:00 & 11/03/2026 01:59 & Não & Limite: fim do horário de saída permitido, sem pernoite (01:59) & R\$ 25,00 \\
CT13 & 10/03/2026 20:00 & 11/03/2026 01:59 & Não & Limite: imediatamente antes da condição de pernoite & R\$ 40,00 \\
CT14 & 10/03/2026 20:00 & 11/03/2026 08:00 & Não & Limite: início da condição de pernoite (08:00 do dia seguinte) & R\$ 50,00 \\
CT15 & 10/03/2026 20:00 & 12/03/2026 08:00 & Não & Partição: duas pernoites completas & R\$ 100,00 \\
CT16 & 10/03/2026 08:00 & 10/03/2026 08:21 & Sim & Partição: cliente VIP com tarifa fixa & R\$ 7,50 \\
CT17 & 10/03/2026 20:00 & 11/03/2026 08:00 & Sim & Partição: cliente VIP com pernoite & R\$ 25,00 \\
CT18 & 10/03/2026 08:00 & 10/03/2026 08:10 & Sim & Partição: cliente VIP dentro da cortesia & R\$ 0,00 \\
\bottomrule
\end{longtable}

\section*{Casos de teste: entrada de veículo fora do horário permitido}

Espera-se que o método lance \texttt{IllegalArgumentException} para todos os casos
a seguir.

\begin{longtable}{@{}p{1.4cm}p{2.8cm}p{2.8cm}p{6.2cm}@{}}
\toprule
\textbf{ID} & \textbf{Entrada} & \textbf{Saída} & \textbf{Técnica / limite} \\
\midrule
\endhead
EI01 & 10/03/2026 07:59 & 10/03/2026 09:00 & Limite: último minuto antes da abertura \\
EI02 & 10/03/2026 00:00 & 10/03/2026 09:00 & Partição: horário de fechamento total \\
EI03 & 10/03/2026 03:00 & 10/03/2026 09:00 & Partição: madrugada, dentro do intervalo inválido \\
\bottomrule
\end{longtable}

\section*{Casos de teste: saída de veículo fora do horário permitido}

Espera-se que o método lance \texttt{IllegalArgumentException} para todos os casos
a seguir. A entrada é deslocada para o dia anterior de modo a isolar a validação do
horário de saída da validação de ordem cronológica.

\begin{longtable}{@{}p{1.4cm}p{2.8cm}p{2.8cm}p{6.2cm}@{}}
\toprule
\textbf{ID} & \textbf{Entrada} & \textbf{Saída} & \textbf{Técnica / limite} \\
\midrule
\endhead
ES01 & 10/03/2026 20:00 & 11/03/2026 02:00 & Limite: início do intervalo de saída proibida (02:00) \\
ES02 & 10/03/2026 20:00 & 11/03/2026 07:59 & Limite: fim do intervalo de saída proibida (07:59) \\
ES03 & 10/03/2026 20:00 & 11/03/2026 05:00 & Partição: madrugada, dentro do intervalo proibido \\
\bottomrule
\end{longtable}

\section*{Caso de teste: ordem cronológica inválida}

\begin{longtable}{@{}p{1.4cm}p{2.8cm}p{2.8cm}p{6.2cm}@{}}
\toprule
\textbf{ID} & \textbf{Entrada} & \textbf{Saída} & \textbf{Técnica / limite} \\
\midrule
\endhead
ET01 & 10/03/2026 10:00 & 10/03/2026 10:00 ou 10/03/2026 09:00 & Saída igual ou anterior à entrada deve lançar \texttt{IllegalArgumentException} \\
\bottomrule
\end{longtable}

\end{document}
